
\section{Preliminary Results}
\label{Early-Results}

We conducted an initial evaluation to validate the first step of the SECDA-DSE loop: translating a natural language accelerator specification (see Appendix) into a SECDA-native implementation that can be compiled through HLS. 
The LLM Stack was provided with a prompt describing an element-wise vector multiplication kernel, where two input vectors $X$ and $Y$ of length $L$ are streamed via AXI-Stream interfaces into on-chip buffers. The accelerator follows a load-compute-store model, where the compute stage performs $Z_i = X_i \odot Y_i$ and streams the output vector $Z$ back to memory. 

Overal, SECDA-DSE successfully generated a complete SECDA-native accelerator workspace, including the SystemC accelerator description, build integration, and a software driver.


%***************************************************************************************************

\subsection{High Level Synthesis}

The evaluation of the generated hardware accelerator was performed using Vivado HLS 2019.2~\cite{vivadohls2019_2} targeting Zynq-7000 FPGA (\texttt{xc7z020-clg400 -1} SoC) with a clock constraint of $5.00$\,ns ($200$\,MHz) and an estimated critical path delay of $3.950$\,ns, meeting the specified clock constraint.


%***************************************************************************************************

\subsection{Latency and Throughput}

The overall accelerator latency ranges from $0$ to $2060$ cycles, corresponding to $0$\,ns to $10.300\,\mu$s at $200$\,MHz. The initiation interval (II) also ranges from $0$ to $2060$ cycles, indicating that the design is not fully pipelined at the top level. 
Table~\ref{tab:vecmul-latency} summarizes the latency characteristics of the submodules.


\begin{table}[t]
  \centering
  \caption{HLS latency and initiation interval (II).}
  \label{tab:vecmul-latency}
  \begin{tabular}{lccc}
    \hline
    Module & Latency (cycles) & II (cycles) & Pipeline \\
    \hline
    HW\_MAIN & 3-3 & 3-3 & None \\
    Send     & 7-1030 & 7-1030 & None \\
    Compute  & 13-1036 & 13-1036 & None \\
    Recv     & 8-2059 & 8-2059 & None \\
    \hline
  \end{tabular}
\end{table}


%***************************************************************************************************


\begin{table}[t]
  \centering
  \caption{HLS resource utilization.}
  \label{tab:vecmul-resources}
  \begin{tabular}{lccc}
    \hline
    Resource & Used & Available & Utilization (\%) \\
    \hline
    BRAM\_18K & 6 & 280 & $\sim$2 \\
    DSP48E    & 3 & 220 & $\sim$1 \\
    FF        & 993 & 106{,}400 & $\sim$0 \\
    LUT       & 1113 & 53{,}200 & $\sim$2 \\
    \hline
  \end{tabular}
\end{table}


\subsection{Resource Utilization}

Table~\ref{tab:vecmul-resources} shows the estimated utilization of the FPGA resources reported by HLS. The compute module accounts for the majority of arithmetic resources, utilizing 3 DSP units along with moderate LUT and FF usage. 
The design maps three on-chip buffers (for $X$, $Y$, and $Z$), each implemented using BRAM. 
These preliminary results demonstrate that SECDA-DSE is capable of generating SECDA-compliant accelerator designs from natural language specifications into SECDA-compliant designs that successfully traverse the HLS toolchain. 
Note that the results are based on HLS estimates, and further evaluation and FPGA execution are part of the full SECDA-DSE scope discussed in Section~\ref{Discussions}.


